\section{Manuscript and literature record}
\subsection{The wording that determines the localizer question}
Our manuscript sources are the supplied transcriptions of cotes 105, 134-8 and 133, the specified edition transcription of \emph{Les D\'erivateurs}, and the typed \emph{Pursuing Stacks} text \citep{GrothArchive105,GrothArchive134,GrothArchive133,GrothDerivateurs,GrothPursuingStacks}. Original scans were not independently inspected. The arguments below can force a mathematical reading, but cannot certify uncertain ink.
We retain the following words from the supplied Q3/Q4 transcriptions because the variance and the three tests are material. In cote 105/052, W7 bis begins: \emph{``Soit $[X_a\leftarrow X_0\to X_b]$ diagramme dans Cat, soit $\Phi=(a\to0\leftarrow b)$ et $\mathcal X$ la cat\'egorie fibr\'ee sur $\Phi$ d\'efinie par le diagramme.''} Its conclusion is
\[
 i_a\in W\Longleftrightarrow i'_a\in W,
 \qquad i_b\in W\Longleftrightarrow i'_b\in W.
\]
The following W8 bis says: \emph{``Consid\'erons les cond. a) [$f\mid f_a$] $\in W$, b) $f_0\in W$, c) $f_a$ et $f_b\in W$. Alors deux de ces conditions entra\^inent la troisi\`eme.''} The brackets record uncertainty in the supplied transcription. A.3 proves that the mathematical reading of a) is the total map $f$; reading it as $f_a$ does not give the cube axiom. The uncertain margin on line 46 is not an additional proved hypothesis.

In XII ms.~75--78 the square axiom uses \emph{``deux ouverts de $X$ tels que $U\cup U'=X$''} and the cube axiom says that \emph{``deux parmi ces conditions impliquent la troisi\`eme''}, with tests at the intersection, at both open pieces, and at the total category. Proposition~6 states that the cube axiom implies the square axiom. Theorem~2 then asserts \emph{``Le localiseur fondamental $W_\infty$ satisfait \`a l'axiome W(7).''} A.8 refutes that last assertion, while A.1--A.3 prove the comparisons with the arbitrary-span formulations.

\subsection{Corrections that affect the answers}
The source augmentation for a sieve cover is the explicit functor $q:T_c(U\leftarrow U\cap V\to V)\to X$. The arrow $\alpha$ in the opposite direction in 105/045--046 is also valid: A.1 constructs it by assigning the intersection the tag $0$ and proves $\alpha\dashv q$. These maps must be distinguished, but the manuscript direction is not an error. The first checkpoint's claim that no canonical intersection tag was available is corrected by that explicit adjunction. The fibred convention is handled by the separate natural comparison in A.1.

The final Der~6, \emph{``cart\'esien ssi cocart\'esien''}, forces the stable pointed-envelope description of A2.4 and hence $RT$-fixedness, but does not force the original derivator to be pointed (A2.3). A2.9 strengthens this conclusion: either unrestricted provisional half alone forces $RT$-fixedness. A2.10 proves equivalence of the constant, locally constant, and arbitrary coefficient versions in the pointed case. In general, A2.13 proves locally constant equivalent to arbitrary, while A2.12 separates constant coefficients by an explicit accessible unpointed derivator. The latter satisfies the constant test for all small indexing categories and fails a locally constant test with five objects and seven arrows. A2.14 further proves that even constant coefficients force $RT$-fixedness for infinity-category derivators, using genuine retractive categories over every coefficient; this enhanced statement is not silently extended to all bare derivators. These distinctions replace the conflicting inferences in R3 and R4. The demand in the 105/051 margin passes to all nonempty shifts automatically for the coherently stated variants, as A2.1 proves. The \emph{``Compl\`etement faux''} margin at 105/052 correctly rejects the preceding assertions about HOT, by A2.2. Its replacement suggestion is itself false: A2.8 gives a bicartesian square over a point whose two middle vertices are both noncontractible.

Grothendieck's footnotes 55 and 56 at XVI ms.~117--118 identify a circular proof of the relative localization axiom. A5.1--A5.4 avoid it: saturation follows from the precofibration-category argument, filtered-colimit closure supplies the relative theorem, and the accessible finite-exception example shows that accessibility alone cannot replace it.

The interval recognition claim in 105/020 and 032 fails even for all sectioned integral-acyclic intervals on simplicial sets (B.4). Its exact additional test is B.2, with the slice-completeness form B.7. In the non-total accessible simplicial case, B.10 identifies the answer exactly as the aspherical core; equality with the given localizer is equivalent to properness. B.12 shows why this criterion cannot be extended to arbitrary strict test categories: interval attachments preserve pointwise support. B.13 separately treats the manuscript's componentwise case (b), using the free chaotic $C_2$ nerve and its cone with fixed apex. This obstruction persists for a proper accessible objectwise localizer and all sectioned intervals. B.11 also corrects an overly literal reading of the one-interval margin: one interval depending on $W$ can encode any set of sectioned intervals, although the fixed Lawvere interval cannot encode every $W$. The small-generating-set condition in 105/040 concerns factorization and is not the recognition test. B.8 records the five requested thesis/book cross-references and their differences in hypotheses.

The literal Lemme~3 at 134-8/053 loses the canonical normalization supplied by condition (c) of the preceding theorem; C.5 refutes the lemma read without that condition. C.11 proves the underlying localization equivalence under the corollary's module hypotheses over every commutative ring; C.10 keeps the exact underived-right-adjoint gate visible. C.13 proves the underlying equivalence for arbitrary AB5 categories with the stated projective and interval hypotheses. C.15 instead assumes an exact cocontinuous resolving functor and removes the colimit-exactness requirement. For the raw adjoint itself, C.7 and C.14 give the ordinary and pure-injective bounds; with scalar-copower compatibility the latter covers every countable coefficient ring. C.16--C.17 prove the raw-adjoint conclusion over every commutative ring when some resolving functor has finite evaluation components up to retract and the chosen interval images are pure. C.19 supplies the alternative disk-coefficient coproduct condition. C.18 shows that the full manuscript assumptions do not imply AB4 or AB5; C.20 refutes a general extension from finite representability to arbitrary sums. Neither obstruction is represented as a counterexample to the normalized full theorem. The additional positive hypotheses are not silently inferred from mere $k$-bilinearity.

Cote 133/006--008 concerns Witt groups and universal lambda formulas, not a claimed finite-local-system localizer. D.1 refutes the naive target-coefficient condition as a localizer; D.2--D.10 give the repair, comparisons and schematization statement. This distinction is necessary to avoid attributing the question in Part D to Grothendieck.

\subsection{The source-by-source opening record}
The file \texttt{DER\_LOCALIZERS\_source\_ledger.md} gives one entry for every source named in the brief and every primary source used in the paper or its recorded exploratory arguments. It distinguishes supplied transcriptions, opened primary texts, metadata-only records, and failed full-text retrievals, with exact theorem numbers and the attempts made. The bibliography repeats the opened/not-obtained status for each entry. Per-part reports and the A3/A4 attempt notes preserve the search strings for A.1--A.5, B.1, C.1 and D.1; no missing search string has been reconstructed and presented as a verbatim log.

In the continuing A.4 search we checked the precise scope of the rigid-group nullity argument, semilocalizing coreflectivity, the acyclic-product framework, and the homotopy-excision estimates \citep{A4Stanley,A4CasacubertaGutierrezRosicky,A4AnelBiedermannFinsterJoyal,A4ChacholskiSchererWerndli}. Their additional hypotheses do not settle the fixed-stable-part question. A34.1--A34.2 instead give the exact ordinary-nullification tests that the remaining separating spaces must satisfy. The report records these unresolved gates and all completed computations. Our literature search is evidence for the sources used, not a priority claim. Downloaded third-party full texts are research inputs and are not redistributed in the archive.

Two further size routes were checked against their original texts.
Casacuberta's nonrepresentable homology kernels fail wedge closure
\citep[Theorem~7.1 and Corollary~7.2]{A4CasacubertaOhkawa2020}.
P\"utzst\"uck's countable-spectrum expression concerns smash-homology
kernels \citep[Theorem~1.4]{A4Putzstuck2026}; BCHSZ explicitly distinguishes
the corresponding stable cohomological questions
\citep[Remarks~2.3, 2.9--2.10]{A4BCHSZ2026}.
Neither yields the required bound on an unstable fixed-stable-part fibre.
The nilpotent Postnikov theorem constructs sections from the space,
and does not supply the opposite inverse-limit passage
\citep[Theorem~6.6]{A4CFFS2015}.
